\section{Controlled Ablation Design}
\label{sec:ablation}

Every Phase~5B condition runs on the same 18 cells with nominal output
capacity of 8,192 tokens.  Table~\ref{tab:conditions} gives the frozen
implementations.

% Evidence: phase5b-frozen-design-and-provider.
\begin{table*}[t]
\centering
\begin{tabular}{@{}clllrr@{}}
\toprule
 & Implementation & Workflow & Instructions & Calls & Max/call \\
\midrule
A & M-short-2 & monolithic & concise & 2 & 4096 \\
B & M-unified-2 & unified & enriched & 2 & 4096 \\
C & M-unified-4 & unified staged & same as B & 4 & 2048 \\
D & Specialized-4 & four roles & role-specific & 4 & 2048 \\
\bottomrule
\end{tabular}
\caption{Frozen conditions.  Each has nominal output capacity 8,192 tokens;
the implementations are not a complete factorial decomposition.}
\label{tab:conditions}
\end{table*}


Condition A is the concise two-call reference.  B preserves that call count
and single unified persona while adding explicit responsibilities for
planning, provenance assessment, review, and reporting.  C retains B's
substantive instruction content but externalizes those responsibilities over
four calls in one continuing workflow.  D assigns the stages to four
specialized roles.  Equal nominal output capacity prevents the four-call
conditions from receiving a larger aggregate token ceiling, although actual
context accumulation and message interfaces still differ.

The four-call workflows externalize planning, provenance assessment, review,
and final reporting.  C executes these responsibilities within a unified
identity and continuing state.  D assigns them to specialized roles with
explicit intermediate interfaces.  A and B use two calls, with the second
producing the final decision after an initial reasoning stage.  Thus the
sequence adds instruction specificity before it adds external stages, and adds
specialization only after the staged unified condition exists.

The contrasts have different interpretations.  \(B-A\) most cleanly isolates
instruction enrichment: both use a single persona and two calls, while B
replaces concise instructions with explicit unified responsibilities.
\(C-B\) holds substantive instructions and nominal output capacity fixed, but
changes two calls into four stages, accumulated context, and persisted
intermediate state.  It therefore evaluates the tested unified multi-call
implementation, not a context-free quantity called reasoning depth.
\(D-C\) compares unified and role-specialized four-call workflows, but also
changes persona scoping, context partitioning, interfaces, and message flow;
it is a bundled residual architectural comparison, not a pure persona
intervention.

The design has 72 intended condition--cell units and 216 planned calls,
\(18(2+2+4+4)\), with zero retries and repairs.  All intended units remain in
the denominator.  Two provider-schema failures are retained under the
preregistered fail-closed interpretation rather than repaired or omitted.
Point estimates are paired all-cell mean regret differences; supporting
analyses use completed pairs, scenario-family comparisons, a two-cluster
bootstrap, and leave-one-extreme-cell trimming.  Because there are only two
scenario clusters, intervals are descriptive rather than population-level
significance claims.  Complete implementation and accounting are in
Supplement Sections E and I.

The execution was frozen before analysis: the registry, visible evidence,
prompts, output schemas, provider configuration, seeds, retry policy, and loss
profile were fingerprinted.  All conditions used the same hosted model and
provider policy with thinking disabled.  The system issued no retries and
performed no schema repairs, so a malformed provider output became a terminal
unit outcome rather than a hidden extra opportunity.  These controls support
within-pilot attribution but do not remove provider nondeterminism or make the
workflow factors factorial.

The preregistered contrast is \(C-B\); \(B-A\) and \(D-C\) are supporting
component contrasts.  We report all-cell estimates, which preserve the
complete denominator under fail-closed treatment, and completed-pair
estimates, which expose the effect of terminal failures without replacing the
canonical analysis.  Scenario and trimming analyses diagnose concentration
rather than serving as post hoc alternatives to the registered estimand.

Equal nominal output capacity is not equal realized computation.  Four-call
conditions accumulate different input contexts and have more synchronization
boundaries.  We therefore avoid describing \(C-B\) as a pure
test-time-compute effect.  Likewise, D introduces role schemas and handoffs in
addition to personas.  The contrast labels encode the observed condition
differences, while the prose states which mechanism each comparison can and
cannot isolate.
